\paragraph{A joint action with dynamical charged matter.}
Prescribed currents can be replaced by a dynamical complex scalar on the
same solid. This requires an additional matter action and coupling;
neither is selected by the classical readback protocol. Gauge-invariant
simplicial actions have established precedents
\cite{ChristiansenHalvorsen2012}. The construction below specifies a
potential-dependent scalar trial space on this particular volume and retains the
complete dependence of the scalar interpolation on the gauge potential.

Let \(a\in\mathbb R^{42}\), \(\phi\in\mathbb R^{13}\) and
\(\psi\in\mathbb C^{13}\) be variable coefficients, and supply a
real charge \(e\). On a tetrahedron, write \(a_{ij}=-a_{ji}\) for
the oriented edge coefficient, set \(a_{ii}=0\), and define
\begin{align*}
 A_h&=W_1a,& \phi_h&=W_0\phi,& \chi_h&=W_0\chi,\\
 \theta_i(x)&=\int_{v_i}^{x}A_h
             =\sum_j a_{ij}\lambda_j(x),&
 \Psi_a(x)&=\sum_i\lambda_i(x)e^{ie\theta_i(x)}\psi_i.
\end{align*}
The paths in \(\theta_i\) are straight segments within the tetrahedron.
The coefficients \(a_{ij}\) are real unwrapped integrals; compact link
phases alone do not determine these interpolants.
Let \(D_t\Psi=\partial_t\Psi+ie\phi_h\Psi\),
\(D_x\Psi=\nabla\Psi-ieA_h\Psi\), and
\(E=-\dot a-D\phi\). With supplied \(m^2,g\geq0\), use the
autonomous action
\begin{equation}
 S=\int dt\left\{\frac12E^{\mathsf T}ME-\frac12a^{\mathsf T}Ka
 +\int_\Omega\left[
 |D_t\Psi_a|^2-|D_x\Psi_a|^2-m^2|\Psi_a|^2
                  -\frac g2|\Psi_a|^4\right]dx\right\}.
 \label{eq:whitney-charged-action}
\end{equation}
Every occurrence of \(\Psi_a\), including its time derivative, is
differentiated when varying \(a\). The field-dependent trial space makes
this requirement substantive.

\begin{theorem}[Gauge-covariant charged field and joint evolution]
\label{thm:whitney-charged-matter}
On the fixed nondegenerate cone, \(\Psi_a\) is a continuous
\(H^1(\Omega)\) field with nodal values \(\psi_i\). The action
\eqref{eq:whitney-charged-action} is invariant under every smooth nodal
gauge transformation
\[
 a\longmapsto a+D\chi,\qquad
 \phi\longmapsto\phi-\dot\chi,\qquad
 \psi_i\longmapsto e^{ie\chi_i}\psi_i.
\]
In temporal gauge, its joint real Euler--Lagrange equations for
\(a,\operatorname{Re}\psi,\operatorname{Im}\psi\) have a unique
global classical solution for every finite initial position and velocity.
If all thirteen Gauss equations hold initially, they hold for all time.
Nonzero locally charged, globally neutral compatible initial data exist
whenever \(e\ne0\). These statements concern the full nonlinear
finite-dimensional action and do not assume a prescribed external current.
\end{theorem}

\begin{proof}
Pullback of a Whitney edge form to the straight segment from \(v_i\)
to \(x\) gives the displayed formula for \(\theta_i\). Since
\(\sum_j\lambda_j=1\), a gauge transformation changes it by
\(\chi_h(x)-\chi_i\), and consequently
\(\Psi_{a+D\chi}[e^{ie\chi}\psi]=e^{ie\chi_h}\Psi_a[\psi]\).
Both covariant derivatives acquire this same unit phase. The electric and
magnetic fields are unchanged, proving action invariance.
At a vertex only its own basis function survives. On a common face only
its vertices contribute, and the barycentric functions and tangential
Whitney field have identical traces from either tetrahedron. The scalar
traces therefore agree. Piecewise smoothness on a finite mesh proves
\(H^1\) conformity.

Write \(\psi=u+iv\) and define the Euler expressions with the sign
convention \(\mathcal E_y=\partial L/\partial y-
d(\partial L/\partial\dot y)/dt\), while
\(\mathcal E_\phi=\partial L/\partial\phi\).
Gauge invariance, integrated against an arbitrary compactly supported
\(\chi(t)\), gives the off-shell identity
\begin{equation}
 \dot{\mathcal E}_\phi+D^{\mathsf T}\mathcal E_a
       +e\bigl(u\mathbin{\odot}\mathcal E_v
                         -v\mathbin{\odot}\mathcal E_u\bigr)=0.
 \label{eq:whitney-matter-noether}
\end{equation}
Here \(\odot\) means coordinatewise multiplication. In particular,
\[
 (\mathcal E_\phi)_i=-(D^{\mathsf T}ME)_i+\rho_i,
 \qquad
 \rho_i=2e\int_\Omega\lambda_i
                   \operatorname{Im}(\overline{\Psi_a}D_t\Psi_a)\,dx.
\]
Thus the joint matter and vector equations propagate Gauss's law. Free
variation of every scalar coefficient requires \(\sum_i\rho_i=0\)
on this closed-load convention; a nonzero total charge would require a
different boundary-flux or source prescription. As for the prescribed
source action, conventional physical charge has the opposite load sign.

To prove existence, fix temporal gauge. Write the scalar interpolation as
\(W(a)\psi\). Its time derivative is
\(W(a)\dot\psi+S(a,\psi)\dot a\), with \(S\) the derivative
of \(W\) in the gauge-potential direction. The quadratic velocity
Hessian is
\[
 \delta\dot a^{\mathsf T}M\delta\dot a
 +2\|W(a)\delta\dot\psi+S(a,\psi)\delta\dot a\|_{L^2}^2.
\]
It is positive definite: vanishing first forces
\(\delta\dot a=0\), and then nodal interpolation forces
\(\delta\dot\psi=0\). All coefficients are smooth real functions,
so the Euler equations form a regular smooth first-order system in
positions and velocities and have unique local solutions.

The conserved energy is the sum of the positive electric and magnetic
energies and \(\|\partial_t\Psi_a\|_{L^2}^2+
\|D_x\Psi_a\|_{L^2}^2+m^2\|\Psi_a\|_{L^2}^2+
g\|\Psi_a\|_{L^4}^4/2\). It bounds \(|\dot a|\), hence
keeps \(a\) in a compact set on each finite time interval.
On that set the positive scalar Gram matrix \(W(a)^*W(a)\) has a
uniform positive lower bound. Smoothness and linearity in \(\psi\)
give \(\|S(a,\psi)\dot a\|_{L^2}\leq C|\psi||\dot a|\).
Consequently \(|\dot\psi|\leq C_T(\sqrt H+|\psi|)\), where
\(H\) is the initial energy. Gronwall's inequality bounds \(\psi\)
and \(\dot\psi\) on every finite interval. Smooth Hessian inversion
then excludes finite-time escape, proving global existence. This argument
does not assert a time-uniform bound on gauge potentials or require a
globally smooth quotient by the gauge group.

For an explicit charged initial condition, take \(a_0=0\),
\(\psi_i(0)=1\), and \(\dot\psi_i(0)=if_i\), where the real
nonzero nodal function \(f_h=\sum_i f_i\lambda_i\) has zero
integral. Antisymmetry gives
\(\sum_{ij}\lambda_i\lambda_j\dot a_{ij}=0\), so the
dressing contribution to \(\partial_t\Psi_a\) cancels at these
data for every \(\dot a\). Hence
\(\rho_i=2e\int\lambda_i f_h\) has zero total and is nonzero
by positivity of the scalar mass matrix. Solve
\(D^{\mathsf T}MDz=\rho\) with mean-zero \(z\), and set
\(E_0=Dz\), \(\dot a_0=-Dz\). This satisfies all Gauss equations
without a circular source definition.
On the regular cone an explicit choice is \(f_0=3\) at the centre and
\(f_i=-1\) at all twelve boundary vertices. If
\(V_T=(3+\sqrt5)/6\) is the common tetrahedron volume, then
\(\rho_0=6eV_T\) and \(\rho_i=-eV_T/2\).
Writing \(\beta=e(2+3\varphi)/10\), take
\(z_0=12\beta/13\) and \(z_i=-\beta/13\).
Each radial electric coefficient is \(-\beta\) and every boundary
electric coefficient is zero. The apex gradient norm
\(|\nabla\lambda_0|^2=3/(2+3\varphi)\), simplex incidence and
the exact barycentric moments verify all thirteen Gauss equations.
\end{proof}

The finite holonomy and interpolation gauge algebra is formalized in
\path{Lean/Screen/WhitneyChargedMatter.lean}. Conformity, the variational
identity and global existence are the analytic argument above. In
particular, the action-derived edge force includes the change of
\(\Psi_a\) when \(a\) varies. The Whitney projection of the continuum
scalar current generally omits these terms even when the nodal matter
equations hold. This finite approximation has dynamical
matter and a joint stationary action. Its geometry, scalar species, charge,
mass, interaction coefficient and continuum matter Lagrangian are supplied.
No executed charged-matter history, source-selected matter content, spatial
error bound, interacting quantum completion or Standard Model identification
is asserted. The free radiative Hilbert space above describes a sector;
it is not a quantization of this full interacting action.
